We study the star height of terms of Kleene algebra with complement, that is, regular expressions extended with complement, where terms are compared by their equational theory with respect to languages: two terms are equivalent if they denote the same language under every valuation of the variables by languages.
We show that the star-height hierarchy is infinite in this setting.
For every $h \ge 0$ we give a term with four variables and star height $h+1$ that is not equivalent to any term of star height at most $h$, even if the latter may use arbitrarily many variables.
This contrasts with the standard interpretation, where letters are interpreted as themselves: there, it has recently been announced that every regular language has generalized star height at most three, so the corresponding hierarchy collapses.
The same separation holds over finite relational models, whose equational theory is incomparable with that over languages.
The proof works in a finite relational semantics that captures the equational theory with respect to languages, and in finite cyclic relational models.
On suitable chain-like relational structures parameterized by an integer, the witness terms define periodic patterns whose periods are values of irreducible polynomials of degree growing exponentially in the star height.
Conversely, using parametric Presburger arithmetic, we show that every term of lower star height has eventual periods given by polynomials whose irreducible factors have strictly smaller degree.
A B\'{e}zout argument then separates the two.
